\documentclass[10pt,a4paper]{amsart}
\usepackage[margin=19mm]{geometry}
\usepackage{amsmath,amssymb,mathtools}
\usepackage[scr=rsfs,cal=euler]{mathalfa}
\usepackage{xfrac}
\usepackage[unicode,hidelinks,bookmarks=false]{hyperref}
\newcommand{\ie}{i.e.\ }
\newcommand{\cf}{cf.\ }
\newcommand{\resp}{respectively\ }
\newcommand{\Subsection}{\subsection}
\providecommand{\nobreakdash}{\nobreak\hbox{-}}
\setlength{\emergencystretch}{2em}
\multlinegap0pt
\usepackage{bm}
\usepackage{bbm}
\usepackage{dsfont}
\usepackage{xspace}
\usepackage{cancel}
\usepackage{mathtools}
\usepackage{amsthm}
\makeatletter
\@ifundefined{thm}{\newtheorem{thm}{Theorem}}{}
\@ifundefined{cor}{\newtheorem{cor}[thm]{Corollary}}{}
\makeatother
\newcommand{\Mo}{\overline{\mathcal{M}}} % Max shuffle
\newcommand{\Mw}{\mathcal{M}_w} % Max shuffle
\newcommand{\M}{\mathcal{M}} % Max shuffle
\newcommand{\cyc}[1]{\operatorname{cyc}_{#1}}
\newcommand{\oln}{\operatorname{oln}}
\newcommand{\set}[1]{\left\{ #1 \right\}}
\newcommand{\tn}{\operatorname{tn }} % Termination Number
\newcommand{\tup}[1]{\left( #1 \right)}
\title[Fixed Point Homing Shuffles]{Fixed Point Homing Shuffles}
\author{Jonathan Parlett}
\date{Source: arXiv:2410.22548v3 (CC BY 4.0)}
\begin{document}
\maketitle

\noindent\emph{Source and licence.} Fixed Point Homing Shuffles. Version arXiv:2410.22548v3 (CC BY 4.0), distributed under the Creative Commons Attribution 4.0 International licence, as recorded in the arXiv licence metadata for this exact version and shown on the abstract page https://arxiv.org/abs/2410.22548v3. Full rights notice: licenses/arxiv-2410.22548-CC-BY-4.0.txt.\par\smallskip

\noindent\textit{Editorial scope.} Counted unit: the Thm whose source label is \texttt{mshufTN} in the article (source0.tex lines 739--745, source label \texttt{mshufTN}), delivered together with the article's own complete original proof of it (source0.tex lines 888--949, one \texttt{proof} environment of 62 lines). No mathematics is written, repaired or re-derived: statement and proof are the article's own text line for line. The proof is detached from the statement in the source: the article states the result at lines 739--745 and prints its proof at lines 888--949, and the proof environment's own header names the statement, so the association is the article's own. Required context retained uncounted: mshufSt2 (lines 727--733); max-term-cor (lines 875--881). Every definition, display and label of those lines is retained unchanged. Omitted from this package and not redistributed: 1--726, 734--738, 746--874, 882--887, 950--984. Nothing inside a retained excerpt is omitted or altered: every retained body is delivered exactly as the article's lines are written, so no omission receipt of any kind accompanies this package. Citations: the retained excerpts carry no citation command at all, so no bibliography entry of the article is transcribed and no reference is redistributed. Reference adaptation: none was needed and none was made; every reference inside the delivered unit resolves inside it, so no unresolved reference or citation is produced. Printed slips kept literal: no printed word, symbol, hypothesis or number of the counted statement or of its proof was altered. Layout and class: the article's own class is \texttt{dmtcs-episciences} with packages including geometry, leftindex, comment, todonotes, amsmath, fancyhdr, graphicx, cancel, amsthm, enumerate, hyperref, xcolor, inputenc, subfigure; the wrapper is the lane's pinned \texttt{amsart} editorial wrapper at 10pt on A4, which restates the theorem-like environments the retained text needs and adds the article's own macro definitions that the retained text needs (\texttt{\textbackslash M}, \texttt{\textbackslash Mo}, \texttt{\textbackslash Mw}, \texttt{\textbackslash cyc}, \texttt{\textbackslash oln}, \texttt{\textbackslash set}, \texttt{\textbackslash tn}, \texttt{\textbackslash tup}), with extra wrapper packages (bm, bbm, dsfont, xspace, cancel, mathtools). The delivered numbering is the unit's own. Assets: no retained span contains a figure environment, a graphics-inclusion command, a PDF-page inclusion, a table or a TikZ picture; no image file is copied into this package, the package declares no asset dependency and no image file is redistributed with it. Licence: version arXiv:2410.22548v3 of this article is distributed under the Creative Commons Attribution 4.0 International licence, as recorded in the arXiv licence metadata for this exact version and shown on the abstract page https://arxiv.org/abs/2410.22548v3. A full rights notice is delivered at licenses/arxiv-2410.22548-CC-BY-4.0.txt. Characters: every retained excerpt is pure ASCII, so the delivered engine, XeLaTeX with its default Latin Modern text font, reports no missing character.
\begin{thm}
    \label{mshufSt2}
    All max shuffles have the same termination number: that is,
    \[
    \tn \M = \tn \Mo \qquad \text{ for any max shuffles } \M,\Mo .
    \]
\end{thm}

\begin{cor}
    \label{max-term-cor}
    Let $ \M $ be a max shuffle. 
    \[
        \text{If } w\tup{1} \ge \Mw\tup{ 1 } \text{ then } \Mw^{n-1}\tup{ 1 } = 1 .
    \]
\end{cor}

\begin{thm}
    \label{mshufTN}
    Let $\M $ be a max shuffle. Assume that $n \ge 2$. Then 
    \[
    \tn \M = 2n-3.
    \]
\end{thm}

\begin{proof}[of Theorem \ref{mshufTN}]
Let $w \in S_n$. We shall show that $\M$ terminates after $2n-3$ iterations on $w$. This will yield $\tn \M \le 2n-3$, thus proving one half of the theorem.

The sequence of front cards
\[
\tup{\Mw^0(1),\ \Mw^1(1),\ \Mw^2(1),\ \dots}
\]
cannot increase forever.
Thus there must be a $d \ge 0$ such that $\Mw^d(1) \ge \Mw^{d+1}(1)$.
Consider the smallest such $d$.
Then the sequence of front cards strictly increases for the first $d$ iterations:
\[
k := w(1) = \Mw^0(1) < \Mw^1(1) < \Mw^2(1) < \cdots < \Mw^d(1).
\]
This entails $d \le n-k$, since the longest increasing sequence of integers starting at $k$ and bounded above by $n$ is the sequence $k,k+1,\dots,n$ of length $n-k+1$. 

Now, if $k=1$, then $\M$ terminates after $0$ iterations on $w$, and hence after $2n-3$ iterations too.
Hence, assume $k \ge 2$. Then $d \le n-k \le n-2$. Now $\Mw^d(1) \ge \Mw^{d+1}(1)$ allows us to apply Corollary \ref{max-term-cor} to $\Mw^d$. This yields $\M^{n-1}_{\M^d_w}\tup{1} = 1$. In other words, $\Mw^{d+n-1}\tup{1} = 1$.
Hence, $\M$ terminates after $d+n-1$ iterations on $w$, and thus after $2n-3$ iterations too (since $d \le n-2$ and so $d+n-1 \le 2n-3$).
This proves $\tn \M \le 2n-3$.

It remains to show $\tn \M \ge 2n-3$.
By Theorem \ref{mshufSt2}, it suffices to show this for one given max shuffle $\M$.
Let $w \in S_n$ be the permutation with one-line notation $(2,3,4,\dots,n,1)$. For example, if $n=5$, then $\oln w = (2,3,4,5,1)$.
We may consider the max shuffle $\M$ described explicitly by
\[
\M\tup{w} = \cyc{w\tup{k},k,k'} w, \qquad \text{ where } k = w\tup{1} \text{ and } k' = \max\tup{w\tup{[2,k]}}
\]
(as long as $k \neq 1$; otherwise $\M\tup{w} = w$; we furthermore understand $\cyc{w\tup{k},k,k'}$ to mean $\cyc{w\tup{k},k}$ when $k' = w\tup{k}$).
It is not hard to see that $\M$ is really a max shuffle.
Informally, $\M$ first swaps the front card $k$ with $w(k)$ before swapping $w(k)$ with largest card $k'$.
Let us use the notation $w \overset{\M}{ \mapsto } w'$ as shorthand for $w' = \M(w)$.
Then we may describe how $\M$ acts on $w=\tup{2,3,4,5,1}$ for $n = 5$ as follows. In following computations we place a dot over $k$, a hat over $w(k)$, and a tilde over $k'$.

\begin{align*}
    \tup{\dot{2},\hat{3},4,5,1} &\overset{\M}{ \mapsto } \tup{\dot{3},2,\hat{4},5,1}\\
     &\overset{\M}{ \mapsto } \tup{\dot{4},2,3,\hat{5},1}\\
     &\overset{\M}{ \mapsto } \tup{\dot{5},2,3,\tilde{4},\hat{1}}\\
     &\overset{\M}{ \mapsto } \tup{\dot{4},2,\tilde{3},\hat{1},5}\\
     &\overset{\M}{ \mapsto } \tup{\dot{3},\tilde{2},\hat{1},4,5}\\
     &\overset{\M}{ \mapsto } \tup{\dot{2},\hat{1},3,4,5}\\
     &\overset{\M}{ \mapsto } \tup{\dot{1},2,3,4,5}.
\end{align*}
The process carries over to arbitrary $n$. Each card (except $1$) appears in front once prior to the $n$-th card and once after.
We are using the notation $[a,b]$ for the list $\tup{a,a+1,\ldots,b-1,b}$.
 \begin{align*}
        (\dot{2},\hat{3},4,\dots,n,1) &\overset{\M}{ \mapsto } (\dot{3},2,\hat{4},\dots,n,1)\\
                          &\overset{\M}{ \mapsto } (\dot{4},2,3,\dots,n,1)\\
                          & \cdots \\
                          &\overset{\M}{ \mapsto } (\dot{k},[2,k-1],[\hat{k+1},n],1) \quad\tup{\text{for each } k \in \set{2,3,\ldots,n} \text{ in order} }\\
                          & \cdots \\
                          &\overset{\M}{ \mapsto } (\dot{n},2,3,\dots,\tilde{n-1},\hat{1})\\
                          &\overset{\M}{ \mapsto } (\dot{n-1},2,3,\dots,\tilde{n-2},\hat{1},n)\\
                          &\overset{\M}{ \mapsto } (\dot{n-2},2,3,\dots,\hat{1},n-1,n)\\
                          & \cdots \\
                          &\overset{\M}{ \mapsto } (\dot{k},[2,\tilde{k-1}],\hat{1},[k+1,n]) \quad\tup{\text{for each } k \in \set{n,n-1,\ldots,1} \text{ in order}}\\
                          & \cdots \\
                          &\overset{\M}{ \mapsto } (\dot{1},2,3,\dots,n-1,n).
\end{align*}
Thus $\M$ terminates in exactly $2n-3$ iterations on $w$.
This proves $\tn \M \ge 2n-3$ for our max shuffle $\M$, and therefore for every max shuffle $\M$, as desired.
\end{proof}

\end{document}
